\documentclass[10pt,a4paper]{amsart}
\usepackage[margin=19mm]{geometry}
\usepackage{amsmath,amssymb,mathtools}
\usepackage[scr=rsfs,cal=euler]{mathalfa}
\usepackage{xfrac}
\usepackage[unicode,hidelinks,bookmarks=false]{hyperref}
\newcommand{\ie}{i.e.\ }
\newcommand{\cf}{cf.\ }
\newcommand{\resp}{respectively\ }
\newcommand{\Subsection}{\subsection}
\providecommand{\nobreakdash}{\nobreak\hbox{-}}
\setlength{\emergencystretch}{2em}
\multlinegap0pt
\usepackage{amssymb}
\usepackage{xcolor}
\usepackage{enumerate}
\usepackage{algorithm}\usepackage{algpseudocode}
\newtheorem{lemma}{Lemma}
\newcommand{\CLSI}[0]{C_{\mathsf{LSI}}}
\newcommand{\E}[0]{\mathbb{E}}
\newcommand{\Ent}[0]{\mathsf{Ent}}
\newcommand{\R}[0]{\mathbb{R}}
\title{Mirror Mean-Field Langevin Dynamics}
\author{Anming Gu, Juno Kim}
\date{Source: arXiv:2505.02621v2 (CC BY 4.0)}
\begin{document}
\maketitle

\noindent\emph{Source and licence.} Mirror Mean-Field Langevin Dynamics. Version arXiv:2505.02621v2 (CC BY 4.0), distributed by arXiv under the Creative Commons Attribution 4.0 International licence, http://creativecommons.org/licenses/by/4.0/, as recorded in the arXiv licence metadata for this exact version and shown on the abstract page https://arxiv.org/abs/2505.02621v2. Full licence notice: \texttt{licenses/arxiv-2505.02621-CC-BY-4.0.txt}.\par\smallskip

\noindent\textit{Editorial scope.} Counted unit: the article's lemma lemma (source0.tex lines 591--593). The article's complete original proof of it is delivered with it (source0.tex lines 594--600, one \texttt{proof} environment of 7 lines). No mathematics is written, repaired or re-derived: the statement and the proof are the article's own lines, in the article's own notation, and no retained line is substituted or re-flowed.  No further source-local context is needed: every cross-reference of every retained line resolves inside the two delivered spans.  Nothing was omitted from any retained span, so the package carries no omission record.  No retained line contains a citation, so the wrapper carries no bibliography and no bibliography entry.  No retained line contains a figure environment, an image-inclusion command or drawing code, so no image is delivered and the package declares no asset dependency.  Editorial formatting only, in addition: an AMS \texttt{amsart} preamble that restates the article's own declarations and macros used by the retained lines, namely the environments \texttt{lemma} on separate counters, together with the standard packages those lines need.  The retained environments print in this document as Lemma 1 in order of appearance, whereas the article prints the same items with the numbering of its own full document.  The complete native source of the article as distributed in the arXiv e-print for this exact version is bundled unchanged as \texttt{sources/177605/source0.tex}.
\begin{lemma}[Convolution]
    Suppose $\mu_1, \mu_2$ respectively satisfy mirror LSI with constant $\CLSI^{(1)}, \CLSI^{(2)}$. Then, their convolution $\mu := \mu_1 * \mu_2$ satisfies a mirror LSI with constant $\left(\frac1{\CLSI^{(1)}}+\frac1{\CLSI^{(2)}}\right)^{-1}$.
\end{lemma}

\begin{proof}
    Let $X\sim \mu_1, Y\sim \mu_2$, and let $f:\R^d\to\R$. Using subadditivity of entropy, \begin{align*}
        \Ent_\mu(f^2)&\le \E[f(X+Y)^2\mid X] + \E[f(X+Y)^2\mid Y]\\
        &\le \left(\frac{2}{\CLSI^{(1)}}+\frac{2}{\CLSI^{(2)}}\right)\E[\|\nabla f\|_{[\nabla^2\phi]^{-1}}^2],
    \end{align*}
which concludes the proof.
\end{proof}

\end{document}
